% =============================================================================
% 3.4 Operational Definitions
% =============================================================================

\section{Operational Definitions}
\label{sec:definitions}

\begin{description}[style=nextline]

  \item[Reliability]
  Three measurable properties:
  (1)~\emph{job success rate}---percentage of Dagster runs completing without error at each concurrency level;
  (2)~\emph{failure isolation}---whether one failing job measurably affects other concurrently running jobs.

  \item[System Stability]
  Measured primarily through \emph{execution time variance}---the standard deviation of job execution time across repeated runs at each concurrency level. A system with high variance is operationally unstable even if its mean execution time is acceptable.

  \item[Operational Overhead]
  Three measurable components:
  (1)~\emph{pod scheduling latency}---time from Dagster job submission to pod entering running state;
  (2)~\emph{container startup time}---time from pod running state to first job instruction executing;
  (3)~\emph{net execution time delta}---difference in total end-to-end execution time between VM and Kubernetes for identical jobs at identical concurrency levels.

  \item[Crossover Point]
  Defined through two complementary thresholds:
  \begin{itemize}
    \item \emph{Reliability crossover:} VM success rate drops below 95\%.
    \item \emph{Performance crossover:} Kubernetes total execution time (including overhead) becomes lower than VM execution time under contention.
    \item \emph{Combined crossover:} Both conditions are simultaneously met.
  \end{itemize}

\end{description}
